% !TEX root = ../../../main.tex
% 来源：高中物理甲种本·第三册（电磁·光学·近代物理）
% 章节：交流电 / 交流电的变化规律
% 原始文件：3.tex，行 139-160
% 类型：tikz
% ---
    \begin{tikzpicture}[>=latex, scale=1.5]
    
    \foreach \x in {.25,.75,-.25,-.75}
    {
        \draw[->] (-2,\x)--(2,\x);
    }
    \draw [dashed](0,-1.5)--(0,2.5)node [right]{中性面};
    \draw (0,0)--(-30:1.5); \draw (0,0)--(150:2.2);
\draw[dashed] (120:1.5)--(-60:1.5);
    \draw[rotate=-60, dashed] (-1.5,-.07) rectangle (1.5,.07);
    \draw[rotate=-30] (-1.5,-.07)node[below]{$d$} rectangle (1.5,.07)node[above]{$a$};
    \draw[rotate=-30, fill=white] (-1.5,0) circle (2.5pt) node{\tiny$\cdot$};
    \draw[rotate=-30, fill=white] (1.5,0) circle (2.5pt) node{\tiny$\times$};
    \draw[rotate=-60, fill=white] (-1.5,0) circle (2.5pt) node{\tiny$\cdot$};
    \draw[rotate=-60, fill=white] (1.5,0) circle (2.5pt) node{\tiny$\times$};
\node at (2.5,0){$B$};

\draw [<->] (0,1.3) arc (90:117:1.3) node [right]{$\phi_0$};
\draw [<->] (0,2) arc (90:150:2) ;
\node at (120:2.2){$\omega t+ \phi_0$};

    \end{tikzpicture}
